\chapter{Inequalities}\index{inequality}

Half of olympiad algebra is inequalities --- and almost all of it grows
from one seed: a real square is nonnegative. This chapter builds the
classical toolkit ($x^2 \ge 0$, AM-GM, Cauchy--Schwarz, Schur, Muirhead),
and just as importantly, the \emph{discipline}: an inequality claim is
incomplete until you name the equality case. After this chapter you can
pick a tool for a target, prove the bound, and exhibit the case where it
is sharp.

\section{What's allowed: the algebra of $<$}\index{inequality!algebra}

From the axioms of an ordered field\index{ordered field}: adding the
same quantity preserves $<$; multiplying by a \emph{positive} quantity
preserves it; multiplying by a \emph{negative} quantity \emph{flips} it;
taking reciprocals of positive quantities flips it; $<$ is transitive.
Every step of an inequality proof must be one of these --- the most
common beginner error is multiplying both sides by an expression whose
sign was never fixed.

\section{Squares are nonnegative}\index{squares, nonnegative}

$(a-b)^2 \ge 0$ with equality iff $a=b$. Expanding gives
$a^2 + b^2 \ge 2ab$: the entire classical apparatus is this fact wearing
cleverer clothes.

\begin{figure}[ht]\centering
\begin{tikzpicture}[scale=0.85]
  % a^2 + b^2 >= 2ab : two squares vs two rectangles, leftover (a-b)^2
  \def\a{3.4}\def\b{1.6}
  % left: a^2 + b^2
  \draw[fill=accent!25] (0,0) rectangle (\a,\a);
  \node at (\a/2,\a/2) {$a^2$};
  \draw[fill=accent!50] (\a+0.3,0) rectangle (\a+0.3+\b,\b);
  \node at (\a+0.3+\b/2,\b/2) {$b^2$};
  \node[below] at (\a/2,-0.15) {$a$};
  \node[below] at (\a+0.3+\b/2,-0.15) {$b$};
  \node[font=\large] at (\a/2-0.2+2.6,-0.9) {area $= a^2+b^2$};
  % right: ab rectangle + ab rectangle + (a-b)^2 hole idea via two ab bars
  \begin{scope}[xshift=7.6cm]
    \draw[fill=good!30] (0,0) rectangle (\a,\b);
    \node at (\a/2,\b/2) {$ab$};
    \draw[fill=good!55] (0,\b+0.2) rectangle (\b,\b+0.2+\a);
    \node[rotate=90] at (\b/2,\b+0.2+\a/2) {$ab$};
    \node[below] at (\a/2,-0.15) {$a$};
    \node[font=\large] at (\a/2+0.4,-0.9) {area $= 2ab$};
  \end{scope}
  % equivalence note
  \node[font=\small,align=center] at (6.4,2.6) {$a^2+b^2-2ab$\\[2pt]$=(a-b)^2 \ge 0$};
  \node[font=\small] at (6.4,4.4) {vs};
\end{tikzpicture}

\caption{$a^2+b^2 \ge 2ab$ as area bookkeeping: the two squares outgrow
the two rectangles by exactly $(a-b)^2$. Notice the equality case is
visible --- the leftover square vanishes exactly when $a = b$.}
\label{fig:sqsum}
\end{figure}

\begin{example}[AM--GM for $n=4$, from scratch]\index{AM--GM!for four}
\textbf{Strategy.} We may not assume AM--GM --- we are proving it. The
only ``free'' inequality is the square. Build up: $n=2$, then bootstrap
four at once, and look at the pattern that generalizes.

\textbf{Solution.} $n=2$: $(x-y)^2\ge0$ gives
$\frac{x+y}{2}\ge\sqrt{xy}$, equality iff $x=y$. For $n=4$:
\[
\frac{a+b+c+d}{4} = \frac{\frac{a+b}{2}+\frac{c+d}{2}}{2}
\ge \frac{\sqrt{ab}+\sqrt{cd}}{2} \ge \sqrt{\sqrt{ab}\sqrt{cd}}
= (abcd)^{1/4},
\]
each step an instance of $n=2$. Equality forces $a=b$, $c=d$, and
$\sqrt{ab}=\sqrt{cd}$, i.e.\ $a=b=c=d$. The shape (``pair up, apply
$n=2$ twice'') is exactly Cauchy's argument: prove $n=2$, prove
$n \Rightarrow 2n$, and \emph{backwards}-fill any $n$ by augmenting with
copies of the mean\index{Cauchy induction}. This is a second specimen
of forward--backward induction (compare \S3).
\end{example}

\section{AM--GM}\index{AM--GM!theorem}

\begin{theorem}[AM--GM]
For $a_1,\dots,a_n \ge 0$,
$\frac{a_1+\cdots+a_n}{n} \ge (a_1\cdots a_n)^{1/n}$, with equality iff
all $a_i$ are equal.
\end{theorem}

Figure~\ref{fig:amgm} shows the classical picture for $n=2$: a chord of
length $a$ then $b$ subtends a semicircle; the radius $(a+b)/2$ always
clears the chord-height $\sqrt{ab}$.
\begin{figure}[ht]\centering
% Figure: semicircle AM-GM for n=2 -- radius vs geometric-mean chord.
% Source for Figure~\ref{fig:amgm}. TikZ only; regenerable by compiling main.tex.
\begin{tikzpicture}[
  lab/.style={font=\footnotesize\sffamily,text=slate},
]
\draw[draw=ink] (0,0) arc (180:0:1.5);
\draw[draw=ink] (-1.5,0) -- (1.5,0);
% split point at x=a vs b: a=1.8? keep simple: a at -0.5 -> segments 1.0,2.0
\draw[accent,very thick] (-0.5,0) -- (-0.5,{sqrt(1.0*2.0)});
\draw[rulegray,dashed] (0,0) -- (0,1.5);
\node[lab] at (-1.0,-0.35) {$a$};
\node[lab] at (0.75,-0.35) {$b$};
\node[lab] at (0.15,1.55) {$\frac{a+b}{2}$};
\node[lab,anchor=west] at (-0.42,0.9) {$\sqrt{ab}$};
\node[lab] at (0,-0.9) {the dashed radius always clears the chord's height};
\end{tikzpicture}

\caption{AM--GM as geometry. Notice the dashed radius $(a+b)/2$ is never
shorter than the perpendicular $\sqrt{ab}$ --- and they coincide exactly
when $a=b$.}
\label{fig:amgm}
\end{figure}

\begin{example}[Minimize $x+\frac1x$ on $x>0$]\index{minimum}
\textbf{Strategy.} Equality cases guide proofs. Guess the symmetric
candidate: $x = 1/x$, i.e.\ $x=1$, value $2$. Now \emph{prove} it ---
which means exhibiting an inequality sharp exactly at $x=1$.

\textbf{Solution.} $x + \frac1x - 2 = \frac{(x-1)^2}{x} \ge 0$ since
$x>0$. Equality iff $x=1$. Two halves: (a) bound $x+1/x \ge 2$ for all
$x>0$; (b) attainability, at $x=1$. Drop either half and the claim is
incomplete --- ``the minimum is 2'' requires both a proof \emph{and} a
witness.
\end{example}

\section{The power-mean chain}\index{power mean}

For $a,b>0$: $\underbrace{\sqrt{\frac{a^2+b^2}{2}}}_{\text{QM}} \ge
\underbrace{\frac{a+b}{2}}_{\text{AM}} \ge \underbrace{\sqrt{ab}}_{\text{GM}}
\ge \underbrace{\frac{2}{1/a+1/b}}_{\text{HM}}$, all equal iff $a=b$.
Figure~\ref{fig:ladder} evaluates the whole ladder at $(2,8)$.
\begin{figure}[ht]\centering
% Figure: power-mean ladder QM >= AM >= GM >= HM on (2,8).
% Source for Figure~\ref{fig:ladder}. TikZ only; regenerable by compiling main.tex.
\begin{tikzpicture}[
  lab/.style={font=\footnotesize\sffamily,text=slate},
  nm/.style={font=\footnotesize\bfseries},
]
% scale: value*0.6 -> QM=sqrt((4+64)/2)=5.83, AM=5, GM=4, HM=32/10=3.2
\draw[rulegray] (0,0) -- (0,4.0);
\foreach \y/\name/\v in {3.50/QM/5.83, 3.00/AM/5, 2.40/GM/4, 1.92/HM/3.2}
{
  \draw[accent,thick] (0,\y) -- (1.4,\y);
  \node[nm,anchor=west] at (1.5,\y) {\name\ $= \v$};
}
\foreach \i in {1,...,3}
  \draw[rulegray] (0.7,0) -- (0.7,4.0);
\node[lab] at (0.7,-0.5) {for $(2,8)$};
\node[lab] at (2.1,-1.1) {four rungs, one theorem --- equality at all four iff $a=b$};
\end{tikzpicture}

\caption{The power-mean ladder on $(2,8)$: QM $5.83 \ge$ AM $5 \ge$ GM $4
\ge$ HM $3.2$. Notice the four rungs collapse to one exactly when $a=b$.}
\label{fig:ladder}
\end{figure}

\section{Cauchy--Schwarz}\index{Cauchy--Schwarz}

\begin{theorem}[Cauchy--Schwarz]
$(\sum a_i^2)(\sum b_i^2) \ge (\sum a_ib_i)^2$, with equality iff the
vectors are proportional.
\end{theorem}

\begin{figure}[ht]\centering
\begin{tikzpicture}[scale=1.15,>=stealth]
  \draw[->] (0,0) -- (3.4,0);
  \draw[->] (0,0) -- (0,2.4);
  \coordinate (B) at (3.0,0);
  \coordinate (A) at (1.9,1.7);
  \draw[->,thick] (0,0) -- (B) node[midway,below] {$\mathbf{b}$};
  \draw[->,thick] (0,0) -- (A) node[midway,above left] {$\mathbf{a}$};
  % projection of a onto b
  \draw[dashed] (A) -- (1.9,0);
  \draw[->,very thick,accent] (0,0) -- (1.9,0) node[midway,below,sloped] {$|\mathbf a|\cos\theta$};
  \draw (0.55,0) arc (0:41.8:0.55) node[right] at (28:0.7) {$\theta$};
  \node[align=left,font=\small] at (4.6,1.3)
    {$\mathbf a\cdot\mathbf b = |\mathbf a||\mathbf b|\cos\theta$\\[2pt]
     $|\mathbf a\cdot\mathbf b| \le |\mathbf a||\mathbf b|$\\[2pt]
     equality $\iff \theta = 0$ (parallel)};
\end{tikzpicture}

\caption{Cauchy--Schwarz is a disguised trigonometric fact:
$\mathbf a\cdot\mathbf b = |\mathbf a||\mathbf b|\cos\theta$ cannot
exceed $|\mathbf a||\mathbf b|$, with equality only when the vectors are
parallel. Notice the discriminant proof manufactures exactly this
$\cos\theta \le 1$ bound.}
\label{fig:cauchyvec}
\end{figure}

Proof: $\sum (a_i t + b_i)^2 \ge 0$ is a quadratic in $t$ with no sign
change, so its discriminant is nonpositive --- that discriminant is
exactly the claim. The \emph{Engel form}\index{Engel form}: for $b_i>0$,
$\sum \frac{a_i^2}{b_i} \ge \frac{(\sum a_i)^2}{\sum b_i}$ --- ``Titu's
lemma'' --- obtained by Cauchy on $(a_i/\sqrt{b_i})^2$ and
$(\sqrt{b_i})^2$.

\begin{example}[Nesbitt's inequality]\index{Nesbitt}
For $a,b,c>0$: $\frac{a}{b+c}+\frac{b}{c+a}+\frac{c}{a+b}\ge\frac32$.

\textbf{Strategy.} The Engel form eats exactly this shape: write each
term as $a^2/(a(b+c))$ so denominators combine.

\textbf{Solution.}
\[
\sum_{cyc}\frac{a^2}{a(b+c)} \ge
\frac{(a+b+c)^2}{2(ab+bc+ca)} \ge \frac32
\iff (a+b+c)^2 \ge 3(ab+bc+ca),
\]
which is $\frac12\sum(a-b)^2\ge0$. Equality iff $a=b=c$.
\end{example}

\section{Equality conditions}\index{equality case}

Every optimization problem is two problems: prove the bound, then find
where it binds. The second half is skipped constantly --- and it is the
half that tells you \emph{which} inequality to use. Apply AM--GM only
where its equality case ($a_i$ all equal) is attainable in the problem's
constraints; otherwise the ``minimum'' you found is unattainable and the
answer is wrong.

\begin{example}[Reading the equality case off the proof]
Show $a^2 + b^2 + c^2 \ge ab + bc + ca$ for all reals, and find exactly
when equality holds.
\par\smallskip
\textbf{Strategy.} Rearrange into a sum of squares; the equality case
then reads itself off --- every square must vanish simultaneously.
\par\smallskip
\begin{solution}
$2(a^2+b^2+c^2) - 2(ab+bc+ca) = (a-b)^2 + (b-c)^2 + (c-a)^2 \ge 0$,
dividing by $2$ gives the claim. Equality needs all three squares zero,
i.e. $a = b = c$. The same two-step structure solves every problem of
this shape: first the bound, then the configuration attaining it ---
and the sum-of-squares form makes the second step mechanical.
\end{solution}
\end{example}

\section{Muirhead and Schur}\index{Muirhead}\index{Schur}

For symmetric sums use the bracket notation\index{symmetric sum!notation}:
$[3,0,0] = \sum_{sym} x^3/2$ conventions vary; we use
$\sum_{cyc} x^3$ style --- spelled out when ambiguity looms. Muirhead:
if exponent vector $\alpha$ \emph{majorizes}\index{majorization} $\beta$
(partial sums of sorted components), then
$\sum_{sym}x^{\alpha} \ge \sum_{sym}x^{\beta}$ for $x_i>0$.
Figure~\ref{fig:major} pictures $[3,0,0]$ majorizing $[2,1,0]$.

\begin{example}[Schur, derived not stated]\index{Schur!derived}
\textbf{Strategy.} Don't invoke Schur as a black box --- manufacture it.
$\sum (a-b)(a-c)(a+b-c)$-type thinking aside, the honest derivation:
expand the nonnegativity of $\sum a(a-b)(a-c)$.

\textbf{Solution.} For $a,b,c\ge0$, claim
$\sum_{cyc} a(a-b)(a-c) \ge 0$: WLOG $a\ge b\ge c$; group the first two
terms: $a(a-b)(a-c)+b(b-c)(b-a) = (a-b)^2(a+b-c) \ge 0$ since $a\ge b$
and $a+b\ge c$; the last term $c(c-a)(c-b) \ge 0$ since both factors are
nonpositive. The full check is exr:i2. Expanding:
\[
\sum_{cyc} a^3 + 3abc \;\ge\; \sum_{sym} a^2 b
\quad\text{(Schur of degree 3),}
\]
equality iff $a=b=c$ or two are equal and the third is $0$. This is the
standard repair for ``AM--GM won't reach'' symmetric cubics.
\end{example}

\begin{figure}[ht]\centering
% Figure: Muirhead majorization -- exponent bars [3,0,0] vs [2,1,0].
% Source for Figure~\ref{fig:major}. TikZ only; regenerable by compiling main.tex.
\begin{tikzpicture}[
  lab/.style={font=\footnotesize\sffamily,text=slate},
  hd/.style={font=\footnotesize\bfseries},
]
% [3,0,0] bars
\node[hd] at (0.6,1.1) {$[3,0,0]$};
\draw[fill=accent!50] (0,0) rectangle (0.4,0.9);
\draw[fill=rulegray!30] (0.5,0) rectangle (0.9,0.02);
\draw[fill=rulegray!30] (1.0,0) rectangle (1.4,0.02);
\node[lab] at (0.7,-0.4) {sym.\ sum $= \sum x^3$};
% [2,1,0] bars
\node[hd] at (4.0,1.1) {$[2,1,0]$};
\draw[fill=accent!50] (3.4,0) rectangle (3.8,0.6);
\draw[fill=rulegray!50] (3.9,0) rectangle (4.3,0.3);
\draw[fill=rulegray!30] (4.4,0) rectangle (4.8,0.02);
\node[lab] at (4.1,-0.4) {sym.\ sum $= \sum_{sym} x^2 y$};
\node[lab] at (2.4,-1.1) {$[3,0,0]$ majorizes $[2,1,0]$:
 $3 \ge 2$ and $3+0 \ge 2+1$ $\Rightarrow$ $\sum x^3 \ge \sum_{sym} x^2 y / 6$};
\end{tikzpicture}

\caption{Majorization: $[3,0,0]$'s bars dominate $[2,1,0]$'s cumulatively.
Notice this is a statement about exponent \emph{distribution}, not values.}
\label{fig:major}
\end{figure}

\section{Strategy selection}\index{strategy!inequalities}

\begin{itemize}
\item \emph{Homogenize}: scale until a free parameter (a sum, a product)
is a fixed value --- legitimate when the inequality is homogeneous.
\item \emph{Normalize}: assume $a+b+c=1$ (or $abc=1$) to eliminate a
constraint before the real work.
\item \emph{Tangent line}: guess where equality hits, fit the tangent of
each summand there, prove $\text{summand} \ge \text{tangent}$ pointwise
(Figure~\ref{fig:tangent}). Sharpness is built in at the touching point.
\end{itemize}
\begin{figure}[ht]\centering
% Figure: tangent-line method -- curve touches its tangent at the equality point.
% Source for Figure~\ref{fig:tangent}. TikZ only; regenerable by compiling main.tex.
\begin{tikzpicture}[
  lab/.style={font=\footnotesize\sffamily,text=slate},
  dot/.style={circle,fill=accent,inner sep=1.8pt},
]
\draw[rulegray,->] (-0.2,0) -- (3.4,0);
\draw[rulegray,->] (0,-0.2) -- (0,2.6);
\draw[ink,thick] (0.25,2.4) parabola bend (1.5,0.9) (3.2,1.9);
\draw[accent,thick,dashed] (0.5,1.45) -- (2.9,0.55);
\node[dot] at (1.5,1.02) {};
\node[lab] at (2.4,1.5) {curve};
\node[lab] at (2.75,0.35) {tangent};
\node[lab] at (1.5,-0.4) {the tangent \emph{touches} at the equality case -- sharpness lives there};
\end{tikzpicture}

\caption{Tangent-line method: the curve sits above its tangent, touching
at the equality point. Notice sharpness is \emph{designed in} by choosing
the tangent at the equality case.}
\label{fig:tangent}
\end{figure}

\section{Summary}

\begin{center}\small
\begin{tabular}{@{}lll@{}}
\toprule
\textbf{Tool} & \textbf{Statement} & \textbf{Equality} \\
\midrule
$x^2\ge0$ & $a^2+b^2\ge2ab$ & $a=b$ \\
AM--GM & $\bar a \ge \sqrt[n]{\prod a_i}$ & all equal \\
Power-mean & QM$\ge$AM$\ge$GM$\ge$HM & all equal \\
Cauchy--Schwarz & $\sum a_i^2\sum b_i^2 \ge (\sum a_ib_i)^2$ & proportional \\
Engel & $\sum a_i^2/b_i \ge (\sum a_i)^2/\sum b_i$ & $a_i \propto b_i$ \\
Schur & $\sum a^3 + 3abc \ge \sum_{sym}a^2b$ & equal, or $a=b,c=0$ \\
Muirhead & majorized exponents $\Rightarrow$ bigger sym sum & all equal \\
\bottomrule
\end{tabular}
\end{center}

\section{Code: numerical sanity checks}\index{ineq\_scan}

An inequality you ``believe'' should survive a Monte-Carlo scan ---
\texttt{ineq\_scan.py} checks Nesbitt and AM-GM on $10\,000$ random
positive triples and grid-searches the $x+1/x$ minimum.
\begin{checkpoint}
Predict: what value will the Nesbitt minimum hover just above, and at
which kind of triple? Then run \texttt{python3 ineq\_scan.py}.
\end{checkpoint}
\begin{lstlisting}
Nesbitt over 10000 random triples: min = 1.5003 >= 1.5
  at (a,b,c) = (3.81, 3.848, 3.951)
AM-GM held on all 10000 triples: True
x + 1/x minimizer on grid: x = 1.0, value = 2.0000
theory: equality at x = 1, minimum value = 2
\end{lstlisting}
The empirical minimum sits just above $1.5$ at a nearly-equal triple ---
exactly the theorem's equality case, found by the experiment rather than
assumed. Sampling never proves, but a violated sample would disprove;
here $10\,000$ samples land where theory says they must.
\texttt{test\_13\_ineq.sh} re-runs and greps the summary lines.

\section{Common mistakes}\index{inequality!mistakes}

\begin{itemize}
\item \emph{Flipping on an unsigned multiplier.} Multiplying by
``$a-b$'' without fixing its sign silently reverses $<$. Fix signs
first --- or add, which never flips.
\item \emph{Bound without witness.} ``Minimum is $2$'' needs both
$x+1/x\ge2$ \emph{and} an $x$ attaining it. Most lost marks are the
missing second half.
\item \emph{AM--GM where equality is impossible.} If the constraint
forces $a\ne b$, AM--GM's equality case is unattainable and the bound
doesn't bind --- pick a tool whose equality case \emph{is} feasible.
\end{itemize}

\section*{Exercises}
\subsection*{Inline}
\begin{exercise}[easy]\label{exr:i1}
Prove $\frac{a+b}{2}\le\sqrt{\frac{a^2+b^2}{2}}$ for $a,b>0$, and state
the equality case, using only $(a-b)^2\ge0$.
\end{exercise}
\begin{solution}
Squaring: equivalent to $(a+b)^2\le2(a^2+b^2)$, i.e.\
$0\le(a-b)^2$. Equality iff $a=b$.
\end{solution}
\begin{exercise}[easy]\label{exr:i2}
Complete the WLOG argument in Schur's inequality: for $a\ge b\ge c\ge0$,
show $\sum_{cyc} a(a-b)(a-c) \ge 0$.
\end{exercise}
\begin{solution}
Group the first two terms:
$a(a-b)(a-c)+b(b-c)(b-a) = (a-b)\bigl[a(a-c)-b(b-c)\bigr] =
(a-b)\bigl[a^2-ac-b^2+bc\bigr] = (a-b)^2(a+b-c) \ge 0$
since $a\ge b$ and $a+b\ge c$. The last term $c(c-a)(c-b) \ge 0$ because
$c-a\le0$ and $c-b\le0$. Total $\ge0$.
\end{solution}
\begin{exercise}[medium]\label{exr:i3}
For $x,y>0$ with $x+y=1$, prove $(1+\frac1x)(1+\frac1y)\ge9$.
\end{exercise}
\begin{solution}
Expand: $1 + \frac{x+y}{xy} + \frac1{xy} = 1 + \frac{2}{xy}$, and
$xy\le\bigl(\frac{x+y}{2}\bigr)^2=\frac14$ by AM--GM, so
$1+\frac2{xy}\ge1+8=9$; equality at $x=y=\frac12$.
\end{solution}
\begin{exercise}[medium]\label{exr:i4}
By Engel's form, prove for $a,b,c>0$:
$\frac{a^2}{b}+\frac{b^2}{c}+\frac{c^2}{a}\ge a+b+c$.
\end{exercise}
\begin{solution}
$\sum a^2/b \ge (a+b+c)^2/(a+b+c) = a+b+c$; equality at $a=b=c$.
\end{solution}
\subsection*{Additional}
\begin{exercise}[easy]\label{exr:a1}
Show $x^2+y^2+z^2\ge xy+yz+zx$; find the equality case.
\end{exercise}
\begin{solution}
Multiply by 2: $(x-y)^2+(y-z)^2+(z-x)^2\ge0$; equality iff $x=y=z$.
\end{solution}
\begin{exercise}[easy]\label{exr:a2}
Prove $\frac{2ab}{a+b}\le\sqrt{ab}$ for $a,b>0$ (HM$\le$GM on two
variables).
\end{exercise}
\begin{solution}
Equivalent to $2\sqrt{ab}\le a+b$, i.e.\ $(\sqrt a-\sqrt b)^2\ge0$.
\end{solution}
\begin{exercise}[easy]\label{exr:a3}
Prove or disprove: if $a\ge b$ then $a^2\ge b^2$ --- and repair the
false version.
\end{exercise}
\begin{solution}
False without signs ($-1\ge-2$). True for $a\ge b\ge0$, or with $|a|\ge|b|$.
\end{solution}
\begin{exercise}[easy]\label{exr:a4}
Minimize $\frac{x^2+2}{\sqrt{x^2+1}}$ on real $x$.
\end{exercise}
\begin{solution}
Let $u=\sqrt{x^2+1}\ge1$: $u+1/u\ge2$ but equality $u=1$ forces $x=0$ ---
it \emph{is} attainable; minimum $2$ at $x=0$.
\end{solution}
\begin{exercise}[medium]\label{exr:a5}
Prove for $a,b,c>0$ with $abc=1$: $a+b+c\ge3$, and extend: minimize
$a+b+c+d$ given $abcd=1$ ---
what does the equality case look like?
\end{exercise}
\begin{solution}
AM--GM directly: $a+b+c\ge3\sqrt[3]{abc}=3$, equality at $a=b=c=1$;
likewise $a+b+c+d\ge4$ with all $=1$.
\end{solution}
\begin{exercise}[medium]\label{exr:a6}
Prove Schur of degree 3 implies: for $a,b,c>0$,
$(a+b-c)(b+c-a)(c+a-b)\le abc$.
\end{exercise}
\begin{solution}
Expand the left minus right; it rearranges exactly to the negation of
Schur's form. (This is the substitution trick behind Ravi.)
\end{solution}
\begin{exercise}[medium]\label{exr:a7}
Minimize $\frac1x+\frac1y$ subject to $x+y=2$, $x,y>0$, and state the
minimizer.
\end{exercise}
\begin{solution}
$\frac{x+y}{xy}=\frac2{xy}\ge\frac2{1}=2$ since $xy\le1$; minimizer
$x=y=1$, value $2$.
\end{solution}
\begin{exercise}[medium]\label{exr:a8}
Verify Muirhead's smallest case on two variables: $[2,0]$ majorizes
$[1,1]$, and show $\sum_{sym}x^2 \ge \sum_{sym}xy$ recovers a familiar
inequality.
\end{exercise}
\begin{solution}
On $(x,y)$: $\sum_{sym}x^2 = 2x^2+2y^2$, $\sum_{sym}xy = 2xy$, so the
claim is $x^2+y^2\ge2xy$ --- the square. Check the majorization:
$2\ge1$ and $2+0=1+1$.
\end{solution}
\begin{exercise}[hard]\label{exr:a9}
Prove the tangent-line method on one summand: show
$\frac{x^3}{1+x^2}\ge x-\frac12$ for $x>0$, equality at $x=1$, and
assemble a proof that for $a+b+c=3$, $a,b,c>0$,
$\sum \frac{a^3}{1+a^2}\ge\frac32$.
\end{exercise}
\begin{solution}
$\frac{x^3}{1+x^2} - \bigl(x-\frac12\bigr) =
\frac{x^3 - (x-\frac12)(1+x^2)}{1+x^2} =
\frac{\frac12(x-1)^2}{1+x^2} \ge 0$ on expanding the numerator:
$x^3 - x - x^3 + \frac12 + \frac12x^2 = \frac12(x^2-2x+1)$. Equality iff
$x=1$. Sum over $a,b,c$ with $a+b+c=3$:
$\sum \frac{a^3}{1+a^2} \ge \sum\bigl(a-\frac12\bigr) = 3-\frac32 =
\frac32$, equality at $a=b=c=1$ --- the tangent was fitted exactly
there.
\end{solution}
\begin{exercise}[hard]\label{exr:a10}
Implement-or-derive: extend \texttt{ineq\_scan.py} to scan the
Schur inequality $\sum a^3+3abc\ge\sum_{sym}a^2b$ on $10^4$ random
triples, and print the minimum of LHS$-$RHS.
\end{exercise}
\begin{solution}
One loop; the scan's minimum should approach $0$ at near-equal triples,
matching the equality case --- a falsification would show up as a
strictly negative minimum.
\end{solution}
\begin{exercise}[hard]\label{exr:a11}
Prove Cauchy's backward step: if AM--GM holds for $2n$ it holds for
$n$, by taking the $n$ numbers plus $n$ copies of their arithmetic mean.
\end{exercise}
\begin{solution}
Let $m=\bar a$. Apply the $2n$ case to
$(a_1,\dots,a_n,m,\dots,m)$: mean is still $m$, geometric mean
$\sqrt[2n]{\prod a_i \cdot m^n}\le m$; squaring out $m^n$ gives
$\prod a_i \le m^n$.
\end{solution}
\begin{exercise}[easy]\label{exr:a12}
For $x,y>0$ with $xy=1$, find the minimum of $x^2+y^2$ and state the
equality case.
\end{exercise}
\begin{solution}
$x^2+y^2 \ge 2xy = 2$ by $a^2+b^2\ge2ab$; equality at $x=y=1$.
\end{solution}
\begin{exercise}[medium]\label{exr:a13}
Show for positive $a,b,c$: $(a+b+c)\left(\frac1a+\frac1b+\frac1c\right)
\ge 9$ --- once by Cauchy--Schwarz in Engel form and once by AM--GM
twice.
\end{exercise}
\begin{solution}
Engel: $\sum 1/a = \sum 1^2/a \ge (1+1+1)^2/(a+b+c)$. Or AM--GM on
$a+b+c \ge 3(abc)^{1/3}$ and on $1/a+1/b+1/c \ge 3/(abc)^{1/3}$, then
multiply. Equality at $a=b=c$ either way.
\end{solution}
\begin{exercise}[medium]\label{exr:a14}
Minimize $\frac{x^2+2}{\sqrt{x^2+1}}$ over real $x$ --- and find the
trap in the naive AM--GM answer.
\end{exercise}
\begin{solution}
AM--GM on $u+1/u$ for $u=\sqrt{x^2+1}\ge1$ suggests $\min=2$ at
$u=1$, i.e. $x=0$ --- attainable, so the answer really is $2$ at
$x=0$; the point is that the equality case must be checked for
reachability (here it is: $u=1$ occurs), unlike the classic
$u \ge 2$ variants where $u=1$ is forbidden.
\end{solution}
\begin{exercise}[hard]\label{exr:a15}
Prove Nesbitt a second way: write $\frac{a}{b+c}+1 =
\frac{a+b+c}{b+c}$ and conclude by AM--HM on the three denominators.
\end{exercise}
\begin{solution}
$\sum \frac{a}{b+c} = \sum\frac{a+b+c}{b+c} - 3$. By
AM--HM (or the $\sum 1/x \ge 9/\sum x$ exercise above) on
$b+c, c+a, a+b$: $\sum \frac{1}{b+c} \ge
\frac{9}{2(a+b+c)}$, so $\sum\frac{a+b+c}{b+c} \ge \frac{9}{2}$ and the
claim follows with equality at $a=b=c$.
\end{solution}

\section*{Notes}
The AM--GM forward--backward induction is Cauchy's original argument
\cite{otis}; the power-mean ladder and equality discipline follow the
Excerpts' chapter on sharpness-vs-tractability \cite{otis}. The
discriminant proof of Cauchy--Schwarz and the Engel form are standard
\cite{engel}; Schur and Muirhead statements follow \cite{otis}. Tangent-line
and normalization strategies are catalogued in \cite{chen-ineq}. For the
``prove the bound + exhibit the case'' two-half structure see
\cite{velleman}. Contest uses: 2024 AIME II P4; 2022 AIME I P15;
2022 AIME I P7; 2024 AMC 12B P13; 2010 AIME I P9; 2025 AIME II P15;
2024 AMC 12B P3; 2025 AIME I P14; 2023 AIME I P2.

\section*{Further reading}
\begin{itemize}
\item \emph{Basics of Olympiad Inequalities} (S.~Riasat) ---
\url{https://web.evanchen.cc/handouts/Ineq/en.pdf} --- free PDF; the
standard olympiad inequalities reference at this level.
\item E.~Chen's inequality handouts --- tangent-line tricks and harder
examples:
\url{https://web.evanchen.cc/handouts/}
\item The AoPS wiki's curated inequalities articles (search ``Olympiad
inequalities'' on):
\url{https://artofproblemsolving.com}
\item A.~Engel, \emph{Problem-Solving Strategies}, inequalities
chapters --- the drill source used throughout the notes.
\item Berkeley CS70 notes ---
\url{https://www.eecs70.org/} --- for the probabilistic and CS-side uses
of the same bounds.
\item R.~Hammack, \emph{Book of Proof} --- for writing the
equality-case half of an inequality proof cleanly:

\url{https://richardhammack.github.io/BookOfProof/}
\end{itemize}

\begin{rulebox}[Watch instead]
\begin{itemize}
\item \textbf{mathophilia --- ``Proving the AM-GM Inequality with ZERO
Calculus''}, 9:22,
\url{https://www.youtube.com/watch?v=bA01HNj2mDw} --- AM--GM for
$n=2$, then induction and Cauchy's forward--backward induction:
the AM--GM sections. (Verified 2026-10-11.)
\item \textbf{Khan Academy --- Linear Algebra playlist} ---
\url{https://www.youtube.com/playlist?list=PLFD0EB975BA0CC1E0}
--- 
contains Sal's Cauchy--Schwarz lectures: the Cauchy--Schwarz section.
(Verified 2026-10-11.)
\end{itemize}
\end{rulebox}

\chapternotes
